\hypertarget{appendix-e38976cbbf86}{}
\section*{Mesh topology and homology}
\section{Mesh topology and homology}\label{appendix-e38976cbbf86-mesh-topology-and-homology}

Mesh \textbf{geometry} concerns positions, lengths, angles and curvature. Mesh \textbf{topology} concerns connectivity: which vertices belong to edges, which edges bound faces, which faces meet, where boundaries exist, and how many connected/loop/void structures the complex contains.

Two meshes can have very different coordinates and the same topology. Two triangles can overlap geometrically without being topologically adjacent if they do not share vertex/edge identity.

\subsection{Simplicial-chain view}\label{appendix-e38976cbbf86-simplicial-chain-view}

For a triangle mesh, vertices are 0-simplices, edges are 1-simplices and faces are 2-simplices. A chain is a formal linear combination of simplices.

Boundary operators satisfy

\[
\partial_k\partial_{k+1}=0,
\]

the algebraic statement that ``the boundary of a boundary is zero.''

Cycles are chains whose boundary vanishes. Boundaries are cycles created as the boundary of a higher-dimensional chain. Homology quotients cycles by boundaries:

\[
H_k=Z_k/B_k.
\]

The Betti number \(\beta_k=\dim H_k\) counts independent homology classes over the chosen coefficient field (for example, \(\mathbb{Z}_2\)); \(Z_k\) is the kernel of the \(k\)-boundary map and \(B_k\) is the image of the \((k+1)\)-boundary map.

For a connected orientable closed surface, \(\beta_0=1\); \(\beta_1\) counts independent one-dimensional homology classes (for genus \(g\), \(\beta_1=2g\)); and \(\beta_2=1\) represents the surface's fundamental two-dimensional class. With a boundary, \(\beta_2\) need not be one.

\subsection{Euler characteristic}\label{appendix-e38976cbbf86-euler-characteristic}

For a finite triangle complex with vertex, edge and face counts \(V,E,F\),

\[
\chi=V-E+F.
\]

For a connected closed orientable surface of genus \(g\),

\[
\chi=2-2g.
\]

But equal Euler characteristic alone does not prove arbitrary complexes are topologically equivalent. It is a checksum, not a complete classifier.

\subsection{Mesh diagnostics}\label{appendix-e38976cbbf86-mesh-diagnostics}

Topology checks used by a geometry pipeline should include:

\begin{itemize}
\tightlist
\item
  boundary-edge count;
\item
  non-manifold edges with more than two incident faces;
\item
  inconsistent face orientation;
\item
  disconnected components;
\item
  vertex fan manifoldness;
\item
  isolated vertices/edges;
\item
  duplicate faces;
\item
  Euler characteristic and, where appropriate, homology.
\end{itemize}

tools/labs/mesh\_topology\_lab.cpp now gives an executable layer for unique edges, half-edge twins, boundary/non-manifold classification and Euler characteristic.

Read the connected explanation in \hyperref[03-visual-computing-i-ch002-from-a-scanned-triangle-soup-to-a-mesh-we-can-trust]{Continuous Notes/University - Worlds, Images and Inference - Graphics, Vision and Photography\#From a scanned triangle soup to a mesh we can trust}.

Related: \hyperref[appendix-faac749c9cdd]{Half-edge topology invariants}, \hyperref[appendix-e4d8ed8d6787]{Persistent homology}, \hyperref[appendix-89db9fe1bc8a]{QEM edge collapse and validity}.
